\documentclass[10pt,a4paper]{amsart}
\usepackage{amsmath,amssymb,amsthm,graphicx}
% For commutative diagrams we use Xy-pic
\usepackage[all]{xy}

% We use 2cell for 2-commutative diagrams.
\xyoption{2cell}
% Theorem environments.
%
\theoremstyle{plain}
\newtheorem{theorem}[subsection]{Theorem}
\newtheorem{proposition}[subsection]{Proposition}
\newtheorem{lemma}[subsection]{Lemma}

\theoremstyle{definition}
\newtheorem{definition}[subsection]{Definition}
\newtheorem{example}[subsection]{Example}
\newtheorem{exercise}[subsection]{Exercise}
\newtheorem{situation}[subsection]{Situation}

\theoremstyle{remark}
\newtheorem{remark}[subsection]{Remark}
\newtheorem{remarks}[subsection]{Remarks}

\numberwithin{equation}{subsection}

% Macros
%
\def\lim{\mathop{\mathrm{lim}}\nolimits}
\def\colim{\mathop{\mathrm{colim}}\nolimits}
\def\Spec{\mathop{\mathrm{Spec}}}
\def\Hom{\mathop{\mathrm{Hom}}\nolimits}
\def\Ext{\mathop{\mathrm{Ext}}\nolimits}
\def\SheafHom{\mathop{\mathcal{H}\!\mathit{om}}\nolimits}
\def\SheafExt{\mathop{\mathcal{E}\!\mathit{xt}}\nolimits}
\def\Sch{\mathit{Sch}}
\def\Mor{\mathop{\mathrm{Mor}}\nolimits}
\def\Ob{\mathop{\mathrm{Ob}}\nolimits}
\def\Sh{\mathop{\mathit{Sh}}\nolimits}
\def\NL{\mathop{N\!L}\nolimits}
\def\CH{\mathop{\mathrm{CH}}\nolimits}
\def\proetale{{pro\text{-}\acute{e}tale}}
\def\etale{{\acute{e}tale}}
\def\QCoh{\mathit{QCoh}}
\def\Ker{\mathop{\mathrm{Ker}}}
\def\Im{\mathop{\mathrm{Im}}}
\def\Coker{\mathop{\mathrm{Coker}}}
\def\Coim{\mathop{\mathrm{Coim}}}

% Boxtimes
%
\DeclareMathSymbol{\boxtimes}{\mathbin}{AMSa}{"02}

%
% Macros for moduli stacks/spaces
%
\def\QCohstack{\mathcal{QC}\!\mathit{oh}}
\def\Cohstack{\mathcal{C}\!\mathit{oh}}
\def\Spacesstack{\mathcal{S}\!\mathit{paces}}
\def\Quotfunctor{\mathrm{Quot}}
\def\Hilbfunctor{\mathrm{Hilb}}
\def\Curvesstack{\mathcal{C}\!\mathit{urves}}
\def\Polarizedstack{\mathcal{P}\!\mathit{olarized}}
\def\Complexesstack{\mathcal{C}\!\mathit{omplexes}}
% \Pic is the operator that assigns to X its picard group, usage \Pic(X)
% \Picardstack_{X/B} denotes the Picard stack of X over B
% \Picardfunctor_{X/B} denotes the Picard functor of X over B
\def\Pic{\mathop{\mathrm{Pic}}\nolimits}
\def\Picardstack{\mathcal{P}\!\mathit{ic}}
\def\Picardfunctor{\mathrm{Pic}}
\def\Deformationcategory{\mathcal{D}\!\mathit{ef}}
\usepackage{cleveref}
\crefname{originalreference}{original reference}{original references}
\Crefname{originalreference}{Original reference}{Original references}
\makeatletter
\newlabel{lemma-category-is-colimit}{{09YU}{source}{}{sourcepointer-1}{}}
\newlabel{lemma-category-is-colimit@cref}{{[originalreference][0][]09YU}{[1][0][]source}{}{}{}}
\makeatother
\crefname{originalreference}{original reference}{original references}
\Crefname{originalreference}{Original reference}{Original references}
\makeatletter
\newlabel{lemma-colimit-affine-sites}{{0EYL}{source}{}{sourcepointer-2}{}}
\newlabel{lemma-colimit-affine-sites@cref}{{[originalreference][0][]0EYL}{[1][0][]source}{}{}{}}
\makeatother
\title{Finite locally constant sheaves on a limit of quasi-compact quasi-separated schemes form the colimit of the categories on the approximants}
\author{The Stacks Project Authors}
\date{stacks-project:etale-cohomology.tex:local-snapshot-20261011; GNU FDL 1.2 or later}
\begin{document}
\maketitle
\noindent\textit{Editorial note.} All delivered mathematical text is the author's own, copied from the pinned Stacks Project chapter file etale-cohomology.tex; no statement or proof step was rewritten, reordered or removed.\par\medskip
\noindent\textit{Editorial note.} The counted proof is the author's complete argument: the proof is written out in full: the reduction of the three cases to essential surjectivity, the trivialising finite affine etale covering, its approximation on some X\_\{i'\} and the spreading of the isomorphism of constructible sheaves to that stage.\par\medskip
\noindent\textit{Editorial note.} The label lemma-category-loc-cst-is-colimit names the counted statement, and the proof environment immediately following it in the source is delivered verbatim as its proof.\par\medskip
\noindent\textit{Editorial note.} The counted statement is carried by the author's own lemma environment in the chapter source.\par\medskip
\noindent\textit{Editorial note.} The statement and proof are copied from the pinned chapter file; the two cited inputs are kept as external pointers with their canonical Stacks tags.\par\medskip
\noindent\textit{Editorial note.} The preamble is the chapter's own native preamble excerpt already used for this source: the Xy-pic setup (native preamble lines 14-18) and the native theorem declarations with the mathematical macros (native preamble lines 154-217).\par\medskip
\noindent\textit{Allowed standard prerequisites (editorial).}\begin{itemize}
\item The limit presentation X = lim X\_i with affine transition morphisms. Same chapter, the identification of the category of constructible sheaves on X as the colimit of those on the X\_i, which the present lemma sharpens to locally constant sheaves. Supplies full faithfulness and the form of the remaining question, namely that a constructible sheaf becoming finite locally constant on X already does so at some stage.
\item Approximation of an etale covering of the limit. Same chapter, the lemma allowing an etale covering of X with affine members to be approximated by an etale covering of some X\_\{i'\}. Supplies the passage from the covering U of X trivialising F to a covering U\_\{i'\} of X\_\{i'\} with the same index set.
\end{itemize}\medskip
\ifdefined\hypertarget\hypertarget{sourcepointer-1}{}\fi
\noindent\textit{Source reference.} Original reference 09YU: sources/stacks-etale-cohomology/pinned-original-source.tex, source lines 12552--12552. Prerequisite pointer invoked at proof line(s) 12652; it names a standard fact of the same chapter that the counted argument uses, and every step of the counted proof built on it is retained verbatim in the delivered body.\par\medskip
\ifdefined\hypertarget\hypertarget{sourcepointer-2}{}\fi
\noindent\textit{Source reference.} Original reference 0EYL: sources/stacks-etale-cohomology/pinned-original-source.tex, source lines 7336--7336. Prerequisite pointer invoked at proof line(s) 12663; it names a standard fact of the same chapter that the counted argument uses, and every step of the counted proof built on it is retained verbatim in the delivered body.\par\medskip
\begin{lemma}
\label{lemma-category-loc-cst-is-colimit}
Let $X = \lim_{i \in I} X_i$ be a limit of a directed
system of schemes with affine transition morphisms.
We assume that $X_i$ is quasi-compact and quasi-separated
for all $i \in I$.
\begin{enumerate}
\item The category of finite locally constant sheaves on $X_\etale$
is the colimit of the categories of finite locally constant sheaves
on $(X_i)_\etale$.
\item The category of finite locally constant abelian sheaves on $X_\etale$
is the colimit of the categories of finite locally constant abelian sheaves
on $(X_i)_\etale$.
\item Let $\Lambda$ be a Noetherian ring. The category of finite type,
locally constant sheaves of $\Lambda$-modules on $X_\etale$
is the colimit of the categories of finite type, locally constant
sheaves of $\Lambda$-modules on $(X_i)_\etale$.
\end{enumerate}
\end{lemma}

\begin{proof}
By Lemma \ref{lemma-category-is-colimit} the functor in each case
is fully faithful. By the same lemma, all we have to show to finish
the proof in case (1) is the following: given a constructible sheaf
$\mathcal{F}_i$ on $X_i$ whose pullback $\mathcal{F}$ to $X$ is
finite locally constant, there exists an $i' \geq i$ such that the
pullback $\mathcal{F}_{i'}$ of $\mathcal{F}_i$ to $X_{i'}$
is finite locally constant. By assumption there exists an \'etale covering
$\mathcal{U} = \{U_j \to X\}_{j \in J}$
such that $\mathcal{F}|_{U_j} \cong \underline{S_j}$ for some finite set $S_j$.
We may assume $U_j$ is affine for all $j \in J$.
Since $X$ is quasi-compact, we may assume $J$ finite.
By Lemma \ref{lemma-colimit-affine-sites} we can find an $i' \geq i$
and an \'etale covering $\mathcal{U}_{i'} = \{U_{i', j} \to X_{i'}\}_{j \in J}$
whose base change to $X$ is $\mathcal{U}$. Then
$\mathcal{F}_{i'}|_{U_{i', j}}$ and $\underline{S_j}$ are
constructible sheaves on $(U_{i', j})_\etale$ whose pullbacks to
$U_j$ are isomorphic. Hence after increasing $i'$ we get
that $\mathcal{F}_{i'}|_{U_{i', j}}$ and $\underline{S_j}$
are isomorphic. Thus $\mathcal{F}_{i'}$ is finite locally constant.
The proof in cases (2) and (3) is exactly the same.
\end{proof}

\end{document}
